% Ссылка на приложение обязательна в тексте (п. 4.11.4)
% Пример ссылки: (приложение~А)
\chapter{ПРОМПТЫ ДЛЯ ФОРМИРОВАНИЯ ДАТАСЕТА}
\label{app:prompts}
 
% В данном приложении приведены полные тексты промптов, использованных
% при формировании датасета для оценки БЯМ-оценщика
% (раздел~\ref{sec:dataset}). Промпты разделены на две группы:
% генерация вопросов с эталонными ответами и рубриками
% (раздел~\ref{app:prompts:generation}) и генерация ответов студентов
% заданного качества для валидационной выборки
% (раздел~\ref{app:prompts:answers}). Промпт, использованный
% БЯМ-оценщиком при выставлении баллов, приведён в
% разделе~\ref{app:prompts:evaluation}.
 
% Промпты приведены в том виде, в котором они передавались модели,
% с сохранением форматирования и переносов строк. Подстановочные
% переменные обозначены в фигурных скобках: \verb|{topic}|,
% \verb|{question}|, \verb|{target_score}| и аналогичные.
 
\section{Генерация вопросов и эталонных ответов}
\label{app:prompts:generation}
 
Для генерации использовалась пара <<системный промпт~+~пользовательский
промпт>>. Системный промпт задаёт роль модели и фиксирует требования к
формату; пользовательский подставляет тему, тип вопроса и количество
генерируемых единиц.
 
\subsection{Системный промпт}
 
Ты — опытный преподаватель машинного обучения в университете.
  
\vspace{1em}
Твоя задача — составлять качественные учебные вопросы открытого
типа и эталонные ответы к ним.
   
\vspace{1em}
Требования к вопросам:

- Вопрос должен требовать развёрнутого текстового ответа
  (не тест, не код)

- Вопрос должен быть конкретным и однозначным

- Уровень сложности — студент третьего курса технического вуза

- Вопрос на русском языке

- Вопрос должен быть достаточно сложным, не касаться базовых
  понятий и определений (например, что такое линейная регрессия
  или почему эта метрика указывает на переобучение),
  более комплексные и неочевидные вопросы

\vspace{1em}
Требования к эталонному ответу:

- Длина: 3–6 предложений, не больше

- Содержит ключевые понятия и формулы
  (при необходимости в LaTeX-нотации)

- Не содержит лишней воды — только суть

- На русском языке
  
\vspace{1em}
Требования к рубрике:

- Ровно 3 описания для каждого балла

- Для каждой оценки (1 балл, 2 балла, 3 балла) должны быть
  описаны критерии, которые должны быть выполнены для получения
  этой оценки

- Формулировки критериев конкретны: не "правильность ответа",
  а "студент указал, что..."

- На русском языке
  
\vspace{1em}
Отвечай строго в формате JSON без каких-либо пояснений до или
после.
 
\subsection{Пользовательский промпт}
 
Составь {n} разных вопросов открытого типа по теме «{topic}».
  
\vspace{1em}
Тип вопроса: \verb|{type_name}|
  
\vspace{1em}
Описание типа: \verb|{type_description}|
  
\vspace{1em}
Верни JSON-массив из {n} объектов. Каждый объект должен иметь
ровно такую структуру:
  
\vspace{1em}
\begin{verbatim}
{
  "question": "<text of question>",
  "reference_answer": "<text of reference answer>",
  "rubric": [
    {"criterion": "<description of criterion>", "max_score": 1},
    {"criterion": "<description of criterion>", "max_score": 2},
    {"criterion": "<description of criterion>", "max_score": 3}
  ]
}
\end{verbatim}
   
\vspace{1em}
Вопросы должны быть разными — не перефразируй один и тот же
вопрос.

Не добавляй ничего вне JSON-массива.

\section{Генерация ответов студентов разного качества}
\label{app:prompts:answers}
 
Для построения валидационной выборки на каждый вопрос генерировались
ответы, соответствующие баллам 1, 2 и 3 рубрики (ответ на 0 баллов
имитировался пустой строкой и не требовал генерации). Использовалась
аналогичная пара промптов с указанием целевого балла и
соответствующего критерия рубрики.
 
\subsection{Системный промпт}
 
Ты — студент, который отвечает на вопрос по машинному обучению.
  
\vspace{1em}
Твоя задача — сгенерировать ответ ровно на заданный балл по
рубрике.
   
\vspace{1em}
Требования:

- Ответ должен быть на русском языке

- Ответ должен выглядеть как естественный ответ студента,
  а не как комментарий проверяющего

- Ответ должен соответствовать целевому критерию и заслуживать
  ровно целевой балл

- Если целевой балл меньше максимального, не добавляй детали,
  которые явно тянут ответ на более высокий балл

- Не упоминай балл, критерии, рубрику или то, что ответ
  специально сгенерирован

- Не пиши слишком формально и грамотно, не используй сложные
  слова или фразы

- Если целевой балл низкий, то пиши максимально просто и
  коротко
   
\vspace{1em}
Отвечай строго JSON-объектом без пояснений до или после:
\verb|{"answer": "<>"}|
 
\subsection{Пользовательский промпт}
 
Вопрос:

\verb|{question}|
   
\vspace{1em}
Эталонный ответ:

\verb|{reference_answer}|
   
\vspace{1em}
Рубрика:

1 балл: \verb|{criterion_1}|

2 балла: \verb|{criterion_2}|

3 балла: \verb|{criterion_3}|
  
\vspace{1em}
Целевой балл: \verb|{target_score}|

Целевой критерий: \verb|{target_criterion}|
  
\vspace{1em}
Сгенерируй ответ студента, который соответствует целевому
критерию и должен получить ровно \verb|{target_score}| балл(а).
Верни только JSON-объект.
 
\section{Оценка ответов студентов}
\label{app:prompts:evaluation}
 
Промпт, использованный БЯМ-оценщиком при выставлении баллов,
приводится в редакции, выбранной по итогам экспериментов с
конфигурациями оценщика (раздел~\ref{sec:eval_experiments}).
 
% TODO: вставить финальную версию системного и пользовательского
% промпта оценщика после фиксации победившей конфигурации.
